% GENERATED FILE. Edit canonical lessons, then run npm run generate:books.
% canonical-source-hash: fnv1a64:b4215bd54c9adae2
% canonical-lessons: UR-C203-mamun, UR-C203-khala, UR-C203-phuphi, UR-C203-bhatija, UR-C203-pota

\chapter{Maternal uncle, Maternal aunt, Paternal aunt, Nephew, Grandson}
\label{ch:ur-maternal-uncle-maternal-aunt-paternal-aunt-nephew-grandson}
\hlchaptermodality{203}

\begin{chapteropening}
\textbf{By the end of this chapter:} \emph{I can say a maternal uncle, a maternal aunt, a paternal aunt, a nephew and a grandson.}

Complete the last lesson of chapter 203: I can say a maternal uncle, a maternal aunt, a paternal aunt, a nephew and a grandson.
\end{chapteropening}

\section[\texorpdfstring{\ur{ماموں} (māmūṅ)}{ماموں (māmūṅ)}]{\ur{ماموں} (māmūṅ) — a maternal uncle}

\label{lesson:UR-C203-mamun}



\begin{quote}
Before the new one: say the Urdu for a taxi, then the Urdu for a rickshaw.
\end{quote}

\subsection*{You'll want to know: \ur{ماموں}}
\textbf{\ur{ماموں}} — \emph{māmūṅ} — \textquotedblleft{}a maternal uncle\textquotedblright{}.

\begin{grammarlens}[title={What you've built}]
One more word for this chapter.
\end{grammarlens}

\subsection*{Your turn}
\begin{itemize}
\raggedright
  \item \emph{Say it:} \emph{māmūṅ}
  \item \emph{Say it:} \emph{māmūṅ}, once more
  \item \emph{Say it:} say \emph{māmūṅ}
  \item \emph{Recall:} say the Urdu for a taxi, then the Urdu for a rickshaw, then say \emph{māmūṅ} again
\end{itemize}

\begin{tcolorbox}[breakable,colback=teal!4,colframe=teal!35!black,title={Before you move on}]
What does \ur{ماموں} mean? (\textquotedblleft{}A maternal uncle\textquotedblright{}.) Say it once more.
\end{tcolorbox}
\section[\texorpdfstring{\ur{خالہ} (khālā)}{خالہ (khālā)}]{\ur{خالہ} (khālā) — a maternal aunt}

\label{lesson:UR-C203-khala}



\begin{quote}
Before the new one: say the Urdu for a rickshaw, then the Urdu for a maternal uncle.
\end{quote}

\subsection*{You'll want to know: \ur{خالہ}}
\textbf{\ur{خالہ}} — \emph{khālā} — \textquotedblleft{}a maternal aunt\textquotedblright{}.

\begin{grammarlens}[title={What you've built}]
One more word for this chapter.
\end{grammarlens}

\subsection*{Your turn}
\begin{itemize}
\raggedright
  \item \emph{Say it:} \emph{khālā}
  \item \emph{Say it:} \emph{khālā}, once more
  \item \emph{Say it:} say \emph{khālā}
  \item \emph{Recall:} say the Urdu for a rickshaw, then the Urdu for a maternal uncle, then say \emph{khālā} again
\end{itemize}

\begin{tcolorbox}[breakable,colback=teal!4,colframe=teal!35!black,title={Before you move on}]
What does \ur{خالہ} mean? (\textquotedblleft{}A maternal aunt\textquotedblright{}.) Say it once more.
\end{tcolorbox}
\section[\texorpdfstring{\ur{پھوپھی} (phuphī)}{پھوپھی (phuphī)}]{\ur{پھوپھی} (phuphī) — a paternal aunt}

\label{lesson:UR-C203-phuphi}



\begin{quote}
Before the new one: say the Urdu for a maternal uncle, then the Urdu for a maternal aunt.
\end{quote}

\subsection*{You'll want to know: \ur{پھوپھی}}
\textbf{\ur{پھوپھی}} — \emph{phuphī} — \textquotedblleft{}a paternal aunt\textquotedblright{}.

\begin{grammarlens}[title={What you've built}]
One more word for this chapter.
\end{grammarlens}

\subsection*{Your turn}
\begin{itemize}
\raggedright
  \item \emph{Say it:} \emph{phuphī}
  \item \emph{Say it:} \emph{phuphī}, once more
  \item \emph{Say it:} say \emph{phuphī}
  \item \emph{Recall:} say the Urdu for a maternal uncle, then the Urdu for a maternal aunt, then say \emph{phuphī} again
\end{itemize}

\begin{tcolorbox}[breakable,colback=teal!4,colframe=teal!35!black,title={Before you move on}]
What does \ur{پھوپھی} mean? (\textquotedblleft{}A paternal aunt\textquotedblright{}.) Say it once more.
\end{tcolorbox}
\section[\texorpdfstring{\ur{بھتیجا} (bhatījā)}{بھتیجا (bhatījā)}]{\ur{بھتیجا} (bhatījā) — a nephew (brother's son)}

\label{lesson:UR-C203-bhatija}



\begin{quote}
Before the new one: say the Urdu for a maternal aunt, then the Urdu for a paternal aunt.
\end{quote}

\subsection*{You'll want to know: \ur{بھتیجا}}
\textbf{\ur{بھتیجا}} — \emph{bhatījā} — \textquotedblleft{}a nephew (brother's son)\textquotedblright{}.

\begin{grammarlens}[title={What you've built}]
One more word for this chapter.
\end{grammarlens}

\subsection*{Your turn}
\begin{itemize}
\raggedright
  \item \emph{Say it:} \emph{bhatījā}
  \item \emph{Say it:} \emph{bhatījā}, once more
  \item \emph{Say it:} say \emph{bhatījā}
  \item \emph{Recall:} say the Urdu for a maternal aunt, then the Urdu for a paternal aunt, then say \emph{bhatījā} again
\end{itemize}

\begin{tcolorbox}[breakable,colback=teal!4,colframe=teal!35!black,title={Before you move on}]
What does \ur{بھتیجا} mean? (\textquotedblleft{}A nephew (brother's son)\textquotedblright{}.) Say it once more.
\end{tcolorbox}
\section[\texorpdfstring{\ur{پوتا} (potā)}{پوتا (potā)}]{\ur{پوتا} (potā) — a grandson (son's son)}

\label{lesson:UR-C203-pota}



\begin{quote}
Before the new one: say the Urdu for a paternal aunt, then the Urdu for a nephew.
\end{quote}

\subsection*{You'll want to know: \ur{پوتا}}
\textbf{\ur{پوتا}} — \emph{potā} — \textquotedblleft{}a grandson (son's son)\textquotedblright{}.

\begin{grammarlens}[title={What you've built}]
One more word for this chapter.
\end{grammarlens}

\subsection*{Your turn}
\begin{itemize}
\raggedright
  \item \emph{Say it:} \emph{potā}
  \item \emph{Say it:} \emph{potā}, once more
  \item \emph{Say it:} say \emph{potā}
  \item \emph{Recall:} say the Urdu for a paternal aunt, then the Urdu for a nephew, then say \emph{potā} again
\end{itemize}

\begin{tcolorbox}[breakable,colback=teal!4,colframe=teal!35!black,title={Before you move on}]
What does \ur{پوتا} mean? (\textquotedblleft{}A grandson (son's son)\textquotedblright{}.) Say it once more.
\end{tcolorbox}
